\section{财产犯罪 - 抢夺罪}
\begin{enumerate}
  \item 抢夺和抢劫性质相同，程度不同。
  \item \zd{对物暴力} ，对人有一定危害性，所以\zd{一般是紧密占有}。
        \begin{enumerate}
          \item \al{甲商场买鞋，试穿时一溜烟跑掉了}。
          \item \al{甲商场买鞋，售货员怀里抱着鞋，一下子扯走了}。
        \end{enumerate}
  \item 公开、秘密都可以。公开是常见情形，但不是必要条件。
        \begin{enumerate}
          \item \al{昏醉抢劫}。
          \item \al{一刀把头砍掉，拿走手机}。
        \end{enumerate}
  \item \zd{携带凶器抢夺定抢劫}。拟制。
        \begin{enumerate}
          \item 凶器。扩大解释，\al{冻鱼}，\zd{杀伤力 + 畏惧感}。
          \item 携带。不能显示，随时使用可能性。
        \end{enumerate}
  \item 既遂。取得控制财物。
\end{enumerate}
